% =============================================================
% conclusion_generale.tex — Conclusion générale et perspectives
% (~3 pages) — À rédiger après stabilisation du chapitre 5.
% =============================================================
\chapter*{Conclusion générale et perspectives}
\addcontentsline{toc}{chapter}{Conclusion générale et perspectives}
\thispagestyle{plain}

\begin{attentedoc}
Le guide pédagogique impose ces éléments explicitement : récapitulation de la démarche
annoncée en introduction ; réponse aux problèmes posés au début (\S\ref{sec:problematique}) ;
problèmes rencontrés lors de la réalisation ; apports techniques et autres ; perspectives
d'approfondissement. À rédiger uniquement une fois le chapitre~\ref{ch:tests} stabilisé.
\end{attentedoc}

\todo{Récapitulation de la démarche complète (un paragraphe, chapitre par chapitre).}

\todo{Réponse explicite à la problématique du \S\ref{sec:problematique} --- reformuler la
question posée puis y répondre point par point à partir des résultats du
chapitre~\ref{ch:tests}.}

\todo{Problèmes rencontrés durant le projet --- reprendre et développer la section
\og Difficultés rencontrées et corrections apportées \fg{} du chapitre~\ref{ch:tests}, sous
un angle plus personnel (apprentissage, ajustement de méthode).}

\todo{Apports techniques et personnels --- compétences acquises, ce que le travail a
apporté à MENAL et à l'étudiant, en cohérence avec les objectifs du \S\ref{sec:objectifs}.}

\todo{Limites assumées du projet --- scalabilité non visée, environnement de laboratoire,
incertitude des mesures --- à énoncer sans en minimiser la portée.}

\todo{Perspectives --- ouverture sur des axes de prolongement réellement pertinents pour
MENAL (ex. extension du modèle de détection, durcissement supplémentaire du pipeline) ---
ne pas reprendre des perspectives génériques d'un autre rapport (IA générative, informatique
quantique, etc.) sans les justifier par le contexte réel du projet.}
